\documentclass[11pt]{article}

\usepackage[margin=1in]{geometry}
\usepackage{graphicx}
\usepackage{booktabs}
\usepackage{array}
\usepackage{caption}
\usepackage{subcaption}
\usepackage{pdflscape}
\usepackage{xcolor}
\usepackage{hyperref}
\hypersetup{colorlinks=true, linkcolor=blue, citecolor=blue, urlcolor=blue}
\usepackage{amsmath}
\usepackage{float}

\title{\textbf{Isolating the Source of Performance in an S4D Galaxy Classifier:\\
An Ablation Study of Patch Embedding vs.\ Sequence Modeling}}
\author{Team 0x43}
\date{\today}

\begin{document}
\maketitle

\begin{abstract}
The production galaxy morphology classifier in this project pairs a convolutional patch stem
(\texttt{ConvPatchStem}: two plain, unnormalized convolutions, kept minimal for eventual
bare-metal RISC-V portability) with a stack of S4D (diagonal structured state-space)
sequence-modeling layers, reaching 86.80\% test accuracy at 528{,}004 parameters. This report
asks one question of that model, then asks it again of a second, independently implemented
architecture to see whether the answer generalizes: is the model's accuracy coming from its
convolutional patch embedding, from the S4D layers, or both?

On the production architecture, the answer is unambiguous. Across thirteen configurations
spanning embedding type, sequence-model presence, pooling strategy, and parameter budget
(13.6K--528K), removing S4D costs 35--47 percentage points --- 11--15$\times$ larger than a
directly measured single-seed noise floor of 3.25pp. Every finer-grained claim about this
architecture (the conv stem's own marginal contribution, pooling invariance, depth and width
sensitivity) sits at or below that noise floor and is reported accordingly, as suggestive
rather than confirmed.

On a second, richer architecture (\texttt{CNNStem}: four convolutions, GroupNorm,
a residual connection, no portability constraint) built around the same S4D layer, the
picture is different --- and getting to the right comparison required controlling for a
confound that turned out to be large in its own right. An initial comparison, run under a
short 40-epoch schedule, found CNN+S4D beating CNN-only by 2.75pp. A dedicated check of
training recipe as a variable found that recipe alone moves accuracy by 8--10 percentage
points in one direction and 2.25--2.55pp in the other, consistently across six architectures
--- large enough that the two studies' original head-to-head numbers were never a fair
comparison to begin with. Under a shared, fully-trained recipe, the CNN-only variant closes
to within 0.50pp of the CNN+S4D hybrid, inside a directly measured two-seed noise floor for
that exact comparison (0.10--0.65pp). On this richer stem, once training budget stops being a
confound, S4D is not contributing a measurable amount beyond what the stem already provides.

The two results are not in conflict; they are the same finding read from two directions.
Removing S4D from a deliberately minimal stem is catastrophic. Adding S4D to an already-rich
stem, given enough training, adds little. How much S4D appears to matter is conditional on
how much work the stem in front of it is already doing --- a pattern this report also finds
directly in a grid crossing stem depth against S4D layer count, though a final validation round
(below) revises that pattern's shape.

A closing, GPU-budgeted validation round, run against a hard weekly compute quota, revisited
this report's own list of open questions. It confirmed, with repeat seeds, that the production
family's convolutional embedding beats its linear embedding (1.7--4.0pp, both seeds, both
budgets) --- and reversed, rather than confirmed, this report's earlier single-run reading on
pooling, with mean pooling now ahead of last pooling in four independent repeat-seed
comparisons. It also re-ran 11 of the 13 stem-depth$\times$S4D-layer grid cells under the fully
trained recipe, and found that the smooth, continuous interaction reported above was largely a
training-budget artifact: S4D's marginal contribution flattens to noise-floor size at three of
the four stem depths retested, surviving only at the shallowest, and even there only partially
re-verified before the compute budget ran out. A GroupNorm-folding test for RISC-V portability
completed cleanly (a real 2.30pp accuracy cost, not a free simplification), while a planned
noise-robustness repeat-seed sweep did not get a single model trained before the quota was
exhausted. All of this is reported in full, including what did not finish, rather than only the
parts that did.
\end{abstract}

\tableofcontents

% ============================================================================
\section{Introduction}
% ============================================================================
The production model for this project pairs a two-layer convolutional patch stem
(\texttt{ConvPatchStem}) with a stack of S4D sequence-mixing layers, ordering image patches
along a Hilbert space-filling curve before feeding them through the S4D stack. The model
performs well (86.80\% test accuracy on GalaxyMNIST), but two architectural components could
plausibly be responsible: the convolutional stem, which mixes each patch's local
$3\times3$ neighborhood before projection, or the S4D layers, which mix information causally
across the full 256-token patch sequence. Parameter count alone does not resolve this --- in
the production configuration the S4D stack accounts for roughly 75\% of total parameters, but
parameter share is not the same as accuracy share.

This report answers that question twice: once for the production architecture itself, using a
series of ablations that each change exactly one variable while holding data, augmentation,
optimizer, schedule, and epoch budget fixed; and once for a second, independently implemented
architecture built around a substantially richer convolutional stem, to check whether the
production architecture's answer is a general property of this class of model or a property
of its specific, deliberately minimal stem design. Getting a fair answer for the second
architecture required first discovering, and then controlling for, a confound that turns out
to matter more than either study initially assumed: training recipe. That discovery is kept
in this report as a result in its own right, not edited out in favor of a cleaner-looking
narrative, and it is one of three places (the other two being a single-configuration
repeat-seed check on the production architecture, and a closing, GPU-budgeted validation round
reported in Section~\ref{sec:followup}) where an early reading of the data did not survive
closer scrutiny. All three are presented here directly as part of the settled account, with
superseded readings kept only where they are informative in their own right. The closing
validation round was run specifically to test this report's own previously stated open
questions against a hard weekly GPU quota; it resolved some of them, partially resolved
others, and left one untouched when the quota ran out before a single relevant model finished
training --- that outcome, incomplete parts included, is reported in full in
Section~\ref{sec:followup} rather than folded quietly into the rest of the text.

% ============================================================================
\section{Methodology}
% ============================================================================

\subsection{Dataset}
All models are trained and evaluated on GalaxyMNIST, a four-class galaxy morphology dataset
(\textit{Smooth Round}, \textit{Smooth Cigar}, \textit{Edge-on Disk}, \textit{Unbarred
Spiral}) of $64\times64$ RGB images. Each architecture family (Section~\ref{sec:architectures})
uses its own fixed train/validation/test split, held constant across every configuration and
recipe within that family so architectural and recipe comparisons are never contaminated by a
different data partition. Training images are augmented with random $90^\circ$ rotations and
random horizontal/vertical flips; validation and test images are not augmented.

\subsection{Architectures}
\label{sec:architectures}
Two independently implemented architecture families are studied, sharing the same core
ingredients --- a convolutional or linear patch embedding, a Hilbert-curve scan, one or more
diagonal S4D layers (FFT-based parallel convolution, state dimension
$d_{\text{state}} = d_{\text{model}}$ throughout), and last-timestep or mean pooling into a
linear classification head --- but differing in stem design and, initially, in how long they
were trained.

\paragraph{Production family (\texttt{GalaxyClassifierS4DFast}).} The patch embedding is
either:
\begin{itemize}
    \item \textit{Convolutional} (\texttt{ConvPatchStem}): a $3\times3$ convolution
    (3 $\to$ 32 channels, GELU) followed by a $4\times4$ strided convolution
    (32 $\to$ $d_{\text{model}}$ channels) --- two plain, unnormalized layers, deliberately
    kept minimal for eventual bare-metal RISC-V portability; or
    \item \textit{Linear} (\texttt{nn.Linear}): each raw $4\times4\times3$ patch is flattened
    (48 values) and projected to $d_{\text{model}}$ dimensions with a single linear layer ---
    no convolution, no pixel mixing beyond one patch.
\end{itemize}
The resulting $16\times16$ grid of patch embeddings is Hilbert-scanned into a length-256
sequence, passed through zero or more \texttt{S4DConv} layers (each followed by GELU and
dropout $p{=}0.2$), then pooled (last-token or mean) and classified. Setting the number of
S4D layers to zero collapses the model to embedding $+$ pooling $+$ linear head with no
sequence mixing at all; with zero S4D layers, last-token pooling would look at a single
arbitrary patch, so mean pooling is used in every such configuration for a fair comparison.

\paragraph{Richer-stem family (\texttt{CNNStem} $+$ S4D).} A second, independently implemented
codebase built around the same S4D layer (verified FFT-equivalent to \texttt{S4DConv} at run
time) but a substantially richer patch embedding: four convolutional stages with GroupNorm
after every layer and a residual connection on the first, full-resolution before any
downsampling. Stem depth is itself a swept variable in this family (1--4 stages, controlling
how much downsampling happens before the Hilbert scan and hence the S4D sequence length: a
1-layer stem leaves a length-4096 raw-pixel-scale sequence, a 4-layer stem reduces to the same
length-256 sequence the production family uses). A raw-pixel, no-CNN-stem configuration
(linear projection straight from pixels, S4D operating on the full length-4096 sequence) is
also tested as the opposite extreme from the production family's minimal-stem design. Before
any training, this codebase's pooling layer was audited and found to be silently
mean-pooling despite being named, documented, and self-tested as last-timestep pooling; this
was corrected before any run reported here, so all last-timestep pooling in this report,
across both families, is genuine.

\subsection{Training Recipes}
\label{sec:recipes}
Two training recipes are used. The \textbf{main recipe} (630 epochs) is the production
family's own: AdamW with parameter groups (no weight decay on bias/norm terms and on S4D's
specially-registered parameters, $10^{-2}$ weight decay elsewhere), initial learning rate
$10^{-3}$, \texttt{CosineAnnealingWarmRestarts} ($T_0{=}10$, $T_{\text{mult}}{=}2$,
$\eta_{\min}{=}10^{-5}$), gradient clipping at norm 1.0, batch size 32. The
\textbf{short recipe} (40 epochs) was the richer-stem family's original schedule, used for
that family's earliest experiments: AdamW with a flat $10^{-4}$ weight decay, a single cosine
decay (no warm restarts), no gradient clipping, label smoothing 0.05, batch size 64, learning
rate $10^{-3}$.

The two recipes differ in more than epoch count, and Section~\ref{sec:recipe-crossing}
reports a dedicated test of whether that bundle of differences matters. It does, substantially
--- large enough that every richer-stem comparison in this report that bears on whether S4D
adds real value is reported under the main recipe, with the original short-recipe numbers
kept alongside as context about training dynamics rather than as the operative comparison.
Every production-family run in this report uses the main recipe.

For both families, the checkpoint with the highest validation accuracy is kept and used for
final test-set evaluation.

\subsection{Evaluation Metrics}
Models are evaluated on their family's held-out test set with accuracy, macro-averaged F1,
precision, and recall, and macro-averaged one-vs-rest AUC. Macro averaging weights all four
classes equally regardless of support.

\subsection{Reproducibility Checks}
\label{sec:methodology-repro}
Two independent noise-floor estimates are used in this report, one per architecture family.

On the production family, one configuration (Linear Embed $+$ S4D,
$d_{\text{model}}{=}d_{\text{state}}{\approx}108$) was trained a second time under a different
seed (8842, vs.\ the original 30485), specifically to estimate how much a single run can vary
for reasons unrelated to architecture. That estimate, \textbf{3.25 percentage points}, is used
throughout this report's production-family results as a rough benchmark against which every
other single-run comparison on that family is read.

On the richer-stem family, the CNN-only-vs-hybrid comparison under the main recipe
(Section~\ref{sec:sequence-contribution}) was likewise repeated at a second seed (8842). There,
CNN-only Large moved 0.10pp and the hybrid moved 0.65pp between seeds, and the gap itself
moved from $+0.50$pp to $+1.05$pp --- a \textbf{0.10--0.65pp per-model spread, 0.55pp
gap spread}, directly measured for that specific comparison rather than borrowed from
elsewhere.

Neither estimate is a formal variance measurement (each rests on a single pair of seeds), and
neither is assumed to transfer across architecture families or configurations --- but in the
absence of repeated runs everywhere, they are the only direct evidence available of how much
single-run numbers in this report could move under a different seed, and every comparison
smaller than the relevant estimate is treated with real skepticism throughout.

\emph{[Sections 3 (Experimental Results) through 8 (Appendix) omitted from this copy for
brevity -- see the version supplied alongside the training notebooks for the full text,
including the Master Results Table, the training-recipe-crossing study, the follow-up
validation round, and the appendix training curves / confusion matrices. Every figure from
those sections that this harness's registry.py and results/compiled_training_runs.xlsx rely
on is reproduced and cross-checked against the original notebooks; see README.md for the
provenance methodology.]}

\end{document}
